%==============================================================================
%  ACADEMIC PAPER TEMPLATE  (article + online supplement)
%------------------------------------------------------------------------------
%  A journal-style article with an optional online supplement. The preamble
%  provides the title page, custom commands, citation and cross-reference
%  setup, theorem/float configuration, and supplement machinery; the body is a
%  short, self-contained example paper that exercises each feature. Replace the
%  metadata and the example content with your own.
%
%  COMPILE (bibliography uses biblatex + biber):
%      latexmk -pdf paper-template.tex
%  or manually:
%      pdflatex paper-template      (alias: lualatex / xelatex also work)
%      biber    paper-template
%      pdflatex paper-template
%      pdflatex paper-template
%
%  By default this file uses the free Cochineal stack and builds on any current
%  TeX Live / MiKTeX. A commented Minion Pro (commercial) stack is included at
%  the end of the "Font face packages" block for my own non-arXiv builds.
%==============================================================================

\documentclass[
    fontsize=12pt,DIV=10,fleqn,
    letterpaper,oneside,paper=portrait,
]{scrartcl}

%%%% Font face packages -------------------------------------------------------
%
%  Four interchangeable stacks. EXACTLY ONE must be active (uncommented).
%  Stack 1 (Cochineal, free) is the default and builds on any current TeX Live.
%  Stack 4 (Minion Pro, commercial) is my personal stack, left commented out.
%
%  The %%arxiv-font%% / %%not-arxiv-font%% tags let a build script swap the
%  active stack automatically: the %%arxiv-font%% lines (Cochineal, free) stay
%  active for an arXiv upload, the %%not-arxiv-font%% lines (Minion Pro) for
%  everything else. Keep them if you use such a script; otherwise ignore.

% 1. Cochineal + Source Sans Pro + Inconsolata   <-- DEFAULT (free)
%    Cochineal is a refined Crimson: a warm, old-style book serif, close in
%    character to Minion. newtxmath's matching option gives coherent math.
\usepackage[p,osf]{cochineal}                           % serif %%arxiv-font%%
\usepackage[cochineal,vvarbb]{newtxmath}                % math  %%arxiv-font%%
\usepackage[semibold,scaled=0.95]{sourcesanspro}        % sans  %%arxiv-font%%
\usepackage[scaled=1.04,varqu,varl]{inconsolata}        % mono  %%arxiv-font%%

% 2. Libertinus serif + sans + mono   (free)
%    The whole family from one source -- the most internally coherent stack.
%    Math is newtxmath with letterforms matched to Libertine; native Libertinus
%    Math would need LuaLaTeX/unicode-math, so under pdflatex this is the route.
%\usepackage{libertinus}                                 % serif + sans + mono
%\usepackage[libertine,vvarbb]{newtxmath}                % math

% 3. STIX Two (Times-like) + Source Sans Pro + Inconsolata   (free)
%    Times-flavoured text with the widest math-symbol coverage of these stacks.
%\usepackage{stix2}                                      % serif + math
%\usepackage[semibold,scaled=0.95]{sourcesanspro}        % sans
%\usepackage[scaled=1.04,varqu,varl]{inconsolata}        % mono

% 4. Minion Pro + Source Sans Pro + Inconsolata   (commercial; my personal stack)
%    My personal stack; Minion Pro must be installed separately. A build script
%    swaps it in (via the %%not-arxiv-font%% tags) for non-arXiv builds.
%\usepackage{newtxmath}                                  % math  %%not-arxiv-font%%
%\usepackage[fullfamily,swash,lf,mathtabular]{MinionPro} % serif %%not-arxiv-font%%
%\usepackage[semibold,scaled=0.95]{sourcesanspro}        % sans  %%not-arxiv-font%%
%\usepackage[scaled=1.04,varqu,varl]{inconsolata}        % mono  %%not-arxiv-font%%


%%%% Boilerplate packages -----------------------------------------------------

\usepackage{scrlayer-scrpage} %%KOMA%%
\usepackage{mathtools}
% tabfigures forces lining (tabular) digits in equations/lists/line numbers
% even when the text font uses old-style figures. Comment out if not installed.
\usepackage[eqno,enum,lineno]{tabfigures}
\usepackage{setspace}
\usepackage{ragged2e}
\usepackage{booktabs}
\usepackage{ntheorem}
% noibid: CMOS author-date style opposes an ibidem mechanism.
\usepackage[backend=biber,authordate,noibid]{biblatex-chicago}
\usepackage{footmisc} % Must load this after biblatex-chicago.
% One neutral link color avoids the visual clutter of type-coded links.
% Coloring the border (not the text) keeps printed text black.
\usepackage[allbordercolors={0.8 0.8 1}]{hyperref}
% nameinlink makes the whole "Figure 1" a link; noabbrev avoids "fig."
\usepackage[nameinlink,noabbrev]{cleveref}
\usepackage{microtype} % improves optical fit


%%%% Common packages ----------------------------------------------------------

\usepackage[para,flushleft]{threeparttable}
\usepackage{dcolumn}
\usepackage{enumitem}
% Loaded [pass] so it does not disturb the main text's KOMA DIV type-area;
% the supplement then switches to 1in margins via \newgeometry.
\usepackage[pass]{geometry}
\usepackage{graphicx}


%%%% Custom packages ----------------------------------------------------------

\usepackage{tikz}



%%%% Boilerplate configuration ------------------------------------------------

% Fix page-number positioning under line spacing %%KOMA%%
\KOMAoptions{onpsinit={\linespread{1}\selectfont}}

% Match \RaggedRight paragraph indent to the normal indent
\setlength{\RaggedRightParindent}{\parindent}

% Allow math to break across pages
\allowdisplaybreaks[1]

% A "Hypothesis" theorem environment (sans-serif bold header)
\theoremheaderfont{\normalfont\sffamily\bfseries}
\newtheorem{hypothesis}{Hypothesis}
\crefname{hypothesis}{Hypothesis}{Hypotheses}
\Crefname{hypothesis}{Hypothesis}{Hypotheses}

% \citeyearlinked{key}: prints (and links) only the year of a reference
\DeclareCiteCommand{\citeyearlinked}
  {\boolfalse{citetracker}%
   \boolfalse{pagetracker}%
   \usebibmacro{prenote}}
  {\ifciteindex
     {\indexfield{indextitle}}
     {}%
   \printtext[bibhyperref]{\printfield[citeyear]{labelyear}}}
  {\multicitedelim}
  {\usebibmacro{postnote}}

% \textcitedash{key}: like \textcite but en-dashes co-authors' names,
% e.g. "Murphy-Topel standard errors" -> "Murphy--Topel standard errors".
\newcommand{\textcitedash}[1]{%
  {%
  \DeclareDelimFormat{finalnamedelim}{\textendash}%
  \textcite{#1}%
  }%
}

% \symfootnote{...}: symbolic footnotes (*, dagger, ...) counted separately
% from numbered footnotes. Handy for title-page acknowledgements/affiliations.
\newcounter{savefootnote}
\newcounter{symfootnote}
\newcommand{\symfootnote}[1]{%
  \setcounter{savefootnote}{\value{footnote}}%
  \setcounter{footnote}{\value{symfootnote}}%
  \ifnum\value{footnote}>8\setcounter{footnote}{0}\fi%
  \let\oldthefootnote=\thefootnote%
  \renewcommand{\thefootnote}{\fnsymbol{footnote}}%
  \footnote{#1}%
  \let\thefootnote=\oldthefootnote%
  \setcounter{symfootnote}{\value{footnote}}%
  \setcounter{footnote}{\value{savefootnote}}%
}


%%%% Common configuration -----------------------------------------------------

% \mybb{1} -> blackboard-bold digit, for indicator functions like \mybb{1}(d=2).
% Math fonts (newtxmath, mathdesign, ...) ship blackboard *letters* but no
% blackboard *digit*, so this selects the DSSerif font directly (free; the
% dsserif files are in TeX Live). Since \usefont bypasses the main font, it
% renders identically under any of the font stacks above -- nothing to swap
% when you switch to the free stack.
\newcommand{\mybb}[1]{{\text{{\usefont{U}{DSSerif}{m}{n}#1}}}}

% A dcolumn type "d" that aligns numbers on the decimal point, e.g. d{2.3}
\newcolumntype{d}[1]{D{.}{.}{#1}}


%%%% Custom configuration -----------------------------------------------------




%%%% Supplement table of contents ---------------------------------------------

% Suppress TOC entries for the main text (the supplement gets its own TOC).
\AtBeginDocument{%
  \addtocontents{toc}{\protect\setcounter{tocdepth}{-1}}%
}


%%%% Supplement sectioning -----------------------------------------------------
%
%  \suppsection / \suppsubsection / \suppsubsubsection number the online
%  supplement as S.1, S.1.1, S.1.1.1 and add the first two levels to the TOC.
%  \cref/\Cref refer to them as "supplement section S.1", etc.

\newcounter{suppsection}
\crefname{suppsection}{supplement section}{supplement sections}
\Crefname{suppsection}{Supplement section}{Supplement sections}
\renewcommand{\thesuppsection}{\arabic{suppsection}} % \cref number: "1"
\newcommand{\suppsection}[1]{%
    \refstepcounter{suppsection}%
    \section*{S.\thesuppsection\quad #1}% display: "S.1"
    \addcontentsline{toc}{section}{S.\thesuppsection\quad #1}%
}

\newcounter{suppsubsection}[suppsection] % linked to the suppsection counter
\crefname{suppsubsection}{supplement subsection}{supplement subsections}
\Crefname{suppsubsection}{Supplement subsection}{Supplement subsections}
\renewcommand{\thesuppsubsection}{\thesuppsection.\arabic{suppsubsection}} % "1.1"
\newcommand{\suppsubsection}[1]{%
    \refstepcounter{suppsubsection}%
    \subsection*{S.\thesuppsubsection\quad #1}% display: "S.1.1"
    \addcontentsline{toc}{subsection}{S.\thesuppsubsection\quad #1}%
}

\newcounter{suppsubsubsection}[suppsubsection] % linked to suppsubsection
\crefname{suppsubsubsection}{supplement subsubsection}{supplement subsubsections}
\Crefname{suppsubsubsection}{Supplement subsubsection}{Supplement subsubsections}
\renewcommand{\thesuppsubsubsection}{\thesuppsubsection.\arabic{suppsubsubsection}}
\newcommand{\suppsubsubsection}[1]{%
    \refstepcounter{suppsubsubsection}%
    \subsubsection*{S.\thesuppsubsubsection\quad #1}% display: "S.1.1.1"
    % subsubsections are not added to the TOC; add a line here if you want them.
}


%%%% Paths --------------------------------------------------------------------
%
%  Folders searched by \input{...} (tables) and \includegraphics{...} (figures).
%  Relative to this file. Adjust to wherever your generated assets live.
\makeatletter
\def\input@path{{tables/}{figures/}}
\makeatother
\graphicspath{{tables/}{figures/}}

\addbibresource{references.bib}


%%%% Metadata -----------------------------------------------------------------

\newcommand{\papertitle}{Commitment and present bias: Experimental evidence}
\newcommand{\paperauthor}{First M.~Last}
\newcommand{\paperdate}{30 June 2026}


\begin{document}

%%%%%%%%%%%%%%%%%%%%%%%%%
%%%% BEGIN MAIN TEXT %%%%
%%%%%%%%%%%%%%%%%%%%%%%%%

% Title page should never carry a page number.
\thispagestyle{empty}

\vspace*{6pt}

\begin{center}% begin centered title

\textbf{\sffamily\LARGE \papertitle\symfootnote{For helpful comments I thank
seminar audiences and colleagues; \textcite{RoeSmith2019} generously shared
data and code. This study was approved under IRB protocol 0000 and
pre-registered at the AEA RCT Registry (AEARCTR-0000000).}}

\vfill

{\large

\paperauthor\symfootnote{Department of Economics, Example University, 100
University Drive, City, ST 00000; \href{mailto:you@example.edu}{you@example.edu}.}

\vspace{\baselineskip}

\paperdate

\vspace{\baselineskip}

}

% Optional status line; delete if unwanted.
\textit{Working paper.} \href{https://example.com/paper.pdf}{Latest version available here.}

\end{center}% end centered title

\vfill
\vspace{-1\baselineskip}

\paragraph{Abstract}
People routinely plan to act patiently and then favor the immediate reward when
the moment arrives. In a laboratory experiment, participants make a sequence of
intertemporal allocation decisions under one of two conditions. We estimate a
quasi-hyperbolic discounting model and find that the treatment lowers the
present-bias parameter and raises demand for a commitment device. The results
suggest that small changes in choice architecture can meaningfully shift
self-control.

\vfill
\vspace{-1\baselineskip}

{\RaggedRight

\paragraph{\emph{JEL} Codes} C91, D15, D91

\paragraph{Keywords} present bias, commitment, intertemporal choice, laboratory experiment

}

\vfill

\clearpage

% 1.5 line spacing for the body.
\setstretch{1.5}

\section{Introduction}\label{sec:intro}

People often plan to act patiently and then choose the immediate reward when
the moment arrives. We ask whether a simple change in how future rewards are
presented amplifies this present bias, and whether it raises the demand for
commitment devices.

Cite in parentheses with \parencite{Doe2020}, or in text with
\textcite{LeeKimPark2021}. Multiple sources collapse neatly:
\parencite{Doe2020,RoeSmith2019,LeeKimPark2021}. For an author whose name is
already in the sentence, print only the linked year, as
\citeauthor{Doe2020} (\citeyearlinked{Doe2020}) shows. Cross-reference sections,
equations, tables, and figures with \texttt{cleveref}: see \cref{sec:model},
\cref{eq:objective}, \cref{tbl:results}, and \cref{fig:example}.

\section{Model}\label{sec:model}

A decision maker values a stream of consumption using quasi-hyperbolic discount
factors. The indicator macro \verb|\mybb{1}(...)| typesets a blackboard ``1'',
e.g.\ $\mybb{1}(x>0)$. The objective is a displayed, labelled equation:%
\begin{align}
    U_t &= u(x_t) + \beta \sum_{\tau=1} \delta^\tau\, u(x_{t+\tau}).
    \label{eq:objective}
\end{align}%
A second equation referenced together with the first (\cref{eq:objective,eq:foc}):%
\begin{equation}
    \frac{\partial U_t}{\partial x_t} = 0.
    \label{eq:foc}
\end{equation}

\section{Hypotheses}\label{sec:hypotheses}

State each pre-registered prediction in the \texttt{hypothesis} environment so
it is numbered and \texttt{cleveref}-referable (see \cref{hyp:main}):

\begin{hypothesis}\label{hyp:main}
The treatment lowers the present-bias parameter relative to the control:
$\beta_{\text{treat}} < \beta_{\text{control}}$.
\end{hypothesis}

\begin{hypothesis}\label{hyp:secondary}
Participants in the treatment are more likely to take up the commitment device.
\end{hypothesis}

\section{Methodology}\label{sec:method}

A custom-labelled enumeration (e.g.\ for a multi-step estimation procedure):

\begin{enumerate}[label=\textbf{\sffamily Step \arabic*:}, align=left, leftmargin=*]
    \item Estimate $\beta$ and $\delta$ for each participant by maximum
    likelihood from the allocation choices.
    \item Compare the estimated $\beta$ across conditions, reporting
    \textcitedash{RoeSmith2019} standard errors to account for the two-stage
    estimator.
\end{enumerate}

\section{Results}\label{sec:results}

\Cref{tbl:results} and \cref{fig:example} report the main results.

% ---- Table: threeparttable + booktabs + dcolumn (decimal-aligned) ----------
\begin{table}\centering
  \begin{threeparttable}
    \caption{Example results table\label{tbl:results}}
    \begin{tabular}{l d{1.3} d{1.3}}
    \toprule
                          & \multicolumn{1}{c}{Model 1} & \multicolumn{1}{c}{Model 2} \\
    \cmidrule(lr){2-2}\cmidrule(lr){3-3}
    Treatment effect      & 0.281 & 0.895 \\
    \quad standard error  & 0.044 & 0.052 \\
    \midrule
    Observations          & \multicolumn{1}{c}{897} & \multicolumn{1}{c}{897} \\
    \bottomrule
    \end{tabular}
    \begin{tablenotes}\footnotesize
      \emph{Note:} Estimated treatment effects, with standard errors in the row
      below. Numbers in the \texttt{d\{1.3\}} columns align on the decimal
      point; wrap headers and integer cells in \verb|\multicolumn{1}{c}{...}|.
    \end{tablenotes}
  \end{threeparttable}
\end{table}

% ---- Figure: TikZ panel (self-contained, no external files) -
\begin{figure}\centering
  \begin{tikzpicture}
    \draw[thick] (0,0) rectangle (6,3);
    \node at (3,1.5) {figure placeholder};
  \end{tikzpicture}
  \caption{Example figure\label{fig:example}}
\end{figure}

% ---- Dropping in externally generated assets (R / Stata / Python) ----------
% A common pattern: \input a booktabs/threeparttable .tex for tables and
% a pgfplots/TikZ .tex (or \includegraphics a PDF) for figures. The \makebox
% lets a wide table or figure overrun the text block, centered:
%
% \begin{table}[p!]\centering
%   \makebox[\textwidth][c]{\input{table_regression.tex}}
% \end{table}
%
% \begin{figure}[p!]\centering
%   \makebox[\textwidth][c]{%
%     \newcommand{\axiswidth}{5in}\newcommand{\axisheight}{4in}%
%     \input{figure_estimates.tex}}
%   \caption{...\label{fig:estimates}}
% \end{figure}
%
% \includegraphics[width=\textwidth]{figure_histogram.pdf}

\section{Conclusion}\label{sec:conclusion}

A simple change in how future rewards are presented shifts the estimated
present-bias parameter and raises commitment demand. We note the usual caveats
of a single laboratory sample and point to field replication as the natural
next step.

\vspace{1\baselineskip}

\setstretch{1} % single spacing for the reference list
{\RaggedRight \printbibliography}


%%%%%%%%%%%%%%%%%%%%%%%%%%
%%%% BEGIN SUPPLEMENT %%%%
%%%%%%%%%%%%%%%%%%%%%%%%%%
%
%  Everything below is the "online supplement": its own title page, table of
%  contents, S-prefixed sectioning and float numbers, and footer page style.
%  Delete this whole block if your paper has no supplement.

% Restore TOC writing for the supplement.
\addtocontents{toc}{\protect\setcounter{tocdepth}{1}}

% New page; smaller font, 1in margins, single spacing.
\clearpage
\KOMAoptions{fontsize=11pt} %%KOMA%%
% 1in margins for the supplement (overrides the main text's KOMA type area),
% giving a wider text block so wide tables/screenshots scale down less.
\newgeometry{margin=1in}
% ALTERNATIVE: instead of fixed 1in margins, let KOMA recompute a DIV-based
% type area for the supplement. Comment out the \newgeometry line above and
% uncomment the two lines below (this is the original template behavior):
%\KOMAoptions{DIV=14} %%KOMA%%
%\recalctypearea %%KOMA%%
\normalfont %%KOMA%%
\linespread{1}\selectfont %%KOMA%%

% Prefix all tables and figures with "S."
\renewcommand{\thetable}{S.\arabic{table}}
\renewcommand{\thefigure}{S.\arabic{figure}}

% Keep hyperref anchors unique between main text and supplement.
\renewcommand{\theHsection}{S.\arabic{section}}
\renewcommand{\theHsubsection}{S.\arabic{section}.\arabic{subsection}}
\renewcommand{\theHsubsubsection}{S.\arabic{section}.\arabic{subsection}.\arabic{subsubsection}}
\renewcommand{\theHfigure}{S.\arabic{figure}}
\renewcommand{\theHtable}{S.\arabic{table}}

% Supplement page style: centered "Supplement page N" footer. %%KOMA%%
\newpairofpagestyles[scrheadings]{suppstyle}{%
  \clearpairofpagestyles
  \cfoot{Supplement page \pagemark}%
}
\pagestyle{suppstyle}

% Restart page numbering for the supplement.
\setcounter{page}{1}

% Supplement title page.
\begin{center}
\vspace*{\baselineskip}
{\sffamily\bfseries\LARGE \papertitle}\\
\vspace{\baselineskip}
{\sffamily\Large Supplement for online publication}\\
\vspace{2\baselineskip}
{\large \paperauthor}\\
\vspace{\baselineskip}
{\large \paperdate}\\
\vspace{2\baselineskip}
\end{center}

\tableofcontents

\suppsection{First supplement section}\label{sec:supp:first}

Supplement sections are numbered S.1, S.2, \ldots\ and appear in the contents
above. Reference them like \cref{sec:supp:first} or \cref{sec:supp:second}.

\suppsubsection{A subsection}\label{sec:supp:first:sub}

Numbered S.1.1.

\suppsubsubsection{A sub-subsection}

Numbered S.1.1.1 (not added to the contents by default).

\suppsection{Additional results}\label{sec:supp:second}

Supplement floats are numbered with an ``S.'' prefix automatically: see
\cref{tbl:supp-example}.

\begin{table}[h]\centering
  \begin{threeparttable}
    \caption{A supplementary table\label{tbl:supp-example}}
    \begin{tabular}{lcc}
    \toprule
    Row & Column A & Column B \\
    \midrule
    First  & 1 & 2 \\
    Second & 3 & 4 \\
    \bottomrule
    \end{tabular}
    \begin{tablenotes}\footnotesize
      \emph{Note:} An example supplementary table.
    \end{tablenotes}
  \end{threeparttable}
\end{table}

\end{document}
